\chapter{Topic: Template Functions, STL Iterators and Algorithms}

\begin{aside}
Material from the remainder of the term, not covered in the original lecture
notes. Builds on templates (Lecture 16), \code{vector} and \code{map}
(Lectures 16 and 19), and the abstract iterator (Lecture 17).
\end{aside}

\section{Template Functions}
Functions can be templates too:
\begin{lstlisting}[style=cpp]
template <typename T> T min(T x, T y) {
    return x < y ? x : y;
}

int a = 1, b = 2;
int c = min(a, b);          // T = int, deduced from the arguments
double d = min(1.5, 2.5);   // T = double
int e = min<int>(1, 2.5);   // T given explicitly (2.5 converted to int)
\end{lstlisting}
The compiler deduces \code{T} from the argument types. The template only works
for types \code{T} that support everything the body uses (here
\code{operator<} and copying); otherwise instantiation fails to compile.

Recall \code{foreach} from Lecture 17, which required our own \code{AbsIter}.
Templates make it fully general, with no inheritance at all:
\begin{lstlisting}[style=cpp]
template <typename Iter, typename Fn>
void foreach(Iter start, Iter finish, Fn f) {
    while (start != finish) {
        f(*start);
        ++start;
    }
}

int a[] = {1, 2, 3, 4, 5};
foreach(a, a + 5, addFive);    // raw pointers are iterators!
vector<int> v{1, 2, 3};
foreach(v.begin(), v.end(), addFive);
\end{lstlisting}
\code{Iter} can be any type supporting \code{!=}, \code{*} and prefix
\code{++} --- including plain pointers. This is how the STL algorithms work.

\section{Iterator Categories}
Different containers support different iterator operations:
\begin{center}
\begin{tabular}{@{}lp{0.38\textwidth}l@{}}
\toprule
Category & Supports & Example \\
\midrule
input & read \code{*it}, \code{++}, single pass & \code{istream\_iterator} \\
output & write \code{*it = x}, \code{++} & \code{ostream\_iterator} \\
forward & read/write, multi-pass, \code{++} & \code{forward\_list} \\
bidirectional & forward $+$ \code{--} & \code{list}, \code{map}, \code{set} \\
random access & bidirectional $+$ \code{it + n}, \code{it[n]}, \code{<} & \code{vector}, pointers \\
\bottomrule
\end{tabular}
\end{center}
Algorithms state which category they need; e.g.\ \code{sort} needs random access,
so it works on a \code{vector} but not on a \code{list}.

\section{The \code{<algorithm>} Library}
A suite of template functions, most of which work over iterator ranges
\code{[start, finish)} (start included, finish excluded).

\begin{lstlisting}[style=cpp]
#include <algorithm>

// for_each: apply f to every element
template <typename Iter, typename Fn>
Fn for_each(Iter start, Iter finish, Fn f);

// find: iterator to the first element == val, or finish if none
template <typename Iter, typename T>
Iter find(Iter first, Iter last, const T &val) {
    for (; first != last; ++first)
        if (*first == val) return first;
    return last;
}

// count: number of elements == val
// copy: copy [first, last) to the range starting at result
template <typename InIter, typename OutIter>
OutIter copy(InIter first, InIter last, OutIter result) {
    while (first != last) *result++ = *first++;
    return result;
}

// transform: result[i] = f(first[i])
template <typename InIter, typename OutIter, typename Fn>
OutIter transform(InIter first, InIter last, OutIter result, Fn f) {
    while (first != last) *result++ = f(*first++);
    return result;
}
\end{lstlisting}

\begin{important}[\code{copy} does not allocate]
\code{copy} and \code{transform} \emph{assume} the destination has room:
\begin{lstlisting}[style=cpp]
vector<int> v{1, 2, 3, 4, 5, 6, 7};
vector<int> w(4);                                  // 4 zeros: room for 4
copy(v.begin() + 1, v.begin() + 5, w.begin());     // w = {2, 3, 4, 5}
\end{lstlisting}
Copying into an empty vector writes past its end: undefined behaviour.
\end{important}

\section{Function Objects}
The parameter \code{Fn f} can be anything that can be called with \code{()}: a
function, or an object of a class that overloads \code{operator()}.
\begin{lstlisting}[style=cpp]
class Plus1 {
  public:
    int operator()(int n) const { return n + 1; }
};

Plus1 p;
p(4);     // 5: p looks like a function

transform(v.begin(), v.end(), w.begin(), Plus1{});
\end{lstlisting}
Unlike functions, function objects can carry state:
\begin{lstlisting}[style=cpp]
class Plus {
    int m;
  public:
    explicit Plus(int m) : m{m} {}
    int operator()(int n) const { return n + m; }
};
transform(v.begin(), v.end(), w.begin(), Plus{5});

class IncreasingPlus {         // state changes between calls
    int m = 0;
  public:
    int operator()(int n) { return n + (m++); }
};
// v = {1, 1, 1} becomes w = {1, 2, 3}
transform(v.begin(), v.end(), w.begin(), IncreasingPlus{});
\end{lstlisting}

\section{Lambdas}
Writing a class for every little function is tedious. A \textbf{lambda}
expression creates an anonymous function object inline:
\begin{lstlisting}[style=cpp]
// how many ints in v are even?
bool even(int n) { return n % 2 == 0; }
int x = count_if(v.begin(), v.end(), even);

int y = count_if(v.begin(), v.end(), [](int n) { return n % 2 == 0; });
\end{lstlisting}
Syntax: \code{[captures](parameters) \{ body \}}. The capture list says which
local variables the lambda may use:
\begin{lstlisting}[style=cpp]
int m = 5;
transform(v.begin(), v.end(), w.begin(), [m](int n) { return n + m; });  // by value
int total = 0;
for_each(v.begin(), v.end(), [&total](int n) { total += n; });          // by ref
\end{lstlisting}
\code{[=]} captures everything used by value, \code{[\&]} by reference.

\section{Iterator Adaptors}
\code{<iterator>} provides iterators that don't walk over a container:
\begin{lstlisting}[style=cpp]
#include <iterator>

// write to a stream: each assignment *it = x prints x followed by ", "
ostream_iterator<int> out{cout, ", "};
copy(v.begin(), v.end(), out);         // prints 1, 2, 3, 4, 5, 6, 7,

// grow a container: each assignment does push_back
vector<int> w;                         // empty!
copy(v.begin(), v.end(), back_inserter(w));   // fine: w grows as needed
\end{lstlisting}
\code{back\_inserter} fixes the ``\code{copy} doesn't allocate'' problem.

\section{More Containers}
\begin{center}
\begin{tabular}{@{}lp{0.62\textwidth}@{}}
\toprule
Container & Notes \\
\midrule
\code{vector<T>} & contiguous array; $\bigO(1)$ amortized \code{push\_back};
                   $\bigO(n)$ insert/erase in the middle \\
\code{deque<T>} & like vector, plus $\bigO(1)$ \code{push\_front} \\
\code{list<T>} & doubly linked list; $\bigO(1)$ insert/erase anywhere given an
                 iterator; no random access \\
\code{map<K,V>} & sorted by key (balanced BST); $\bigO(\log n)$ operations \\
\code{set<T>} & sorted unique elements \\
\code{unordered\_map<K,V>} & hash table; expected $\bigO(1)$ operations; unordered \\
\bottomrule
\end{tabular}
\end{center}

\section{Example}
Read words from stdin, and print each distinct word with its frequency, sorted:
\begin{lstlisting}[style=cpp]
#include <iostream>
#include <map>
#include <string>
using namespace std;

int main() {
    map<string, int> freq;
    string w;
    while (cin >> w) ++freq[w];
    for (auto &p : freq) cout << p.first << " " << p.second << endl;
}
\end{lstlisting}

\section*{Summary}
\begin{itemize}
  \item Template functions deduce their type parameters from the arguments.
  \item STL algorithms work on iterator ranges \code{[first, last)} and any
        type with the right operations, including pointers.
  \item Iterator categories determine which algorithms apply.
  \item Pass behaviour as functions, function objects (\code{operator()}), or
        lambdas, written
        \code{[captures](params)\{body\}}.
  \item \code{copy}/\code{transform} don't allocate; use \code{back\_inserter}
        or size the destination first; \code{ostream\_iterator} prints.
\end{itemize}
